\documentclass[12pt]{article}
\usepackage[utf8]{inputenc}
\usepackage[T1]{fontenc}
\usepackage[margin=1in]{geometry}
\usepackage{setspace}
\usepackage{xcolor}
\usepackage[colorlinks=true, urlcolor=blue!60!black, linkcolor=blue!60!black]{hyperref}
\setlength{\parindent}{0em}
\setlength{\parskip}{0.6em}
\renewcommand{\baselinestretch}{1.0}
\begin{document}
\begin{center}
{\LARGE \textit{Korematsu v. United States}, 323 U.S. 214 (1944)\par}
\vspace{0.5em}
{\large Case Brief --- Two-Page Edition\par}
\end{center}

\section{Memory Jogger}
The Court upheld a Japanese American citizen's conviction for defying a World War II exclusion order, deferring to the war powers while announcing that racial restrictions are ``immediately suspect''; it was later repudiated as ``gravely wrong the day it was decided'' (\href{https://www.courtlistener.com/opinion/104040/korematsu-v-united-states/}{323 U.S. 214}).

\section{Detailed Case Facts}
Under Executive Order 9066 and a 1942 statute criminalizing violations of military-area restrictions, General DeWitt's Civilian Exclusion Order No.~34 excluded all persons of Japanese ancestry, including citizens, from part of California and directed them to Assembly Centers. Fred Korematsu, a citizen and welder in San Leandro whose loyalty was never questioned, stayed home and was arrested.

\section{Procedural History}
Korematsu was convicted and placed on probation, then sent to the camps. The Ninth Circuit affirmed, and the Supreme Court affirmed. His conviction was vacated on \textit{coram nobis} in 1984, and \textit{Trump v. Hawaii} (2018) repudiated the decision.

\section{Judicial Votes}
Affirmed, 6--3. Black wrote for Stone, Reed, Frankfurter, Douglas, and Rutledge; Frankfurter concurred; Roberts, Murphy, and Jackson dissented separately.

\section{Holding}
Racial restrictions are ``immediately suspect'' and subject to ``the most rigid scrutiny,'' but when Congress and the President act jointly in wartime against grave danger, temporary exclusion of citizens of Japanese ancestry from a threatened area is permissible. The Court decided only exclusion, not detention.

\section{Analysis of Opinions}
Black deferred to the military's judgment that disloyal persons could not be quickly separated from the loyal, severed exclusion from detention, and denied racial hostility. Frankfurter stressed the war power. Roberts saw a ``cleverly devised trap'' of contradictory orders aimed at imprisonment based on ancestry. Murphy showed that the exclusion rested on racial stereotypes rather than evidence. Jackson warned that judicial approval would turn a military expedient into a lasting constitutional principle.

\section{Examples: Future Applications}
\textbf{Same-Side:} ordering citizens of one national origin to holding centers without hearings is unlawful under \textit{Trump v. Hawaii} and \textit{Endo}. \textbf{Opposite-Side:} a statutory entry suspension for foreign nationals was upheld in \textit{Trump v. Hawaii}. \textbf{Fence-Sitter:} heightened screening of citizens who visited a hostile nation, falling heavily on one ethnic community.

\section{Critique}
Scholars across the spectrum place \textit{Korematsu} in the constitutional ``anticanon,'' faulting the Court's abandonment of text for wartime pragmatism and its failure when the government withheld evidence; Harold Hongju Koh warns that \textit{Trump v. Hawaii}'s deference risks resurrecting its ``ghost.''

\section{Key Quotations}
\begin{quote}
``It should be noted, to begin with, that all legal restrictions which curtail the civil rights of a single racial group are immediately suspect. That is not to say that all such restrictions are unconstitutional. It is to say that courts must subject them to the most rigid scrutiny. Pressing public necessity may sometimes justify the existence of such restrictions; racial antagonism never can.''
\end{quote}
\noindent --- Justice Black, Opinion of the Court.

\end{document}
